\documentclass[11pt,a4paper]{article}

\usepackage[utf8]{inputenc}
\usepackage[T1]{fontenc}
\usepackage{lmodern}
\usepackage[english]{babel}
\usepackage{geometry}
\usepackage{graphicx}
\usepackage{float}
\usepackage{caption}
\usepackage{amsmath}
\usepackage{amssymb}
\usepackage{booktabs}
\usepackage{hyperref}
\usepackage{subcaption}
\usepackage{xcolor}

\geometry{margin=1in}
\raggedbottom
\emergencystretch=3em

% Figures live in the repo's runs/ directory; this file sits in report/10/,
% so resolve \includegraphics paths relative to the repository root.
\graphicspath{{../../}}

% Graceful fallback for figures produced by runs that may not be pulled yet.
\newcommand{\figorbox}[2]{%
  \IfFileExists{../../#1}{\includegraphics[width=#2\textwidth]{#1}}{%
    \fbox{\parbox[c][3.5cm][c]{#2\textwidth}{\centering\itshape
    \textcolor{gray}{[figure not yet generated --- run the eval and pull it]}}}}}

\title{Working Notes X \\
       \large \textcolor{red}{PERSONAL WORKING DOCUMENT --- NOT A REPORT FOR THE PROFESSOR,
       NOT PAPER TEXT} \\
       Read-Out-Only Fine-Tuning on Three-Way Splits: Targeted vs.\ Full-Distribution}
\author{Edoardo Paccagnella}
\date{\today}

\begin{document}

\maketitle

\begin{center}
\fbox{\parbox{0.9\textwidth}{\centering\bfseries
This document exists so Edoardo can read Claude's analysis of one specific experiment before
deciding whether/how it goes into the paper (\texttt{idea\_paper\_by\_prof.tex},
\S5.3 ``Generalization to three paths, finetuning??''). It is not written for the professor and
is not paper prose --- do not copy sentences from here into the paper without rewriting them to
the paper's own conventions and evidentiary bar.}}
\end{center}

% =====================================================================
% STATUS: SKELETON. The fine-tuning run (scripts/
% r10_finetune_threeway_splitchains_seed1000.sbatch) has not been launched/pulled yet.
% Fill in every [PENDING] block below once runs/report10/... has real data, using
% plot_finetune_readout_threeway.py to regenerate the three figures locally (no GPU
% needed once the run is pulled).
% =====================================================================

\section{Where this comes from}
\label{sec:intro}

Report~IX, Thread~A.3 (\texttt{report/9/transformer\_for\_graphs\_9.tex}, and the mechanistic
follow-up \texttt{stagewise\_threeway.py}) tested the split-chains-only $n{=}46$ seed-$1000$
checkpoint on every three-way partition of the canvas. Two cells stood out on inspection of the
five-stage similarity-geometry heatmaps:

\begin{itemize}\setlength\itemsep{0.15em}
  \item \textbf{$(15,15,16)$} --- fully balanced. Every within-component block in the final-stage
        geometry is already visibly darker/more positive than the cross-component blocks (it
        \emph{looks} like the model has separated the three components correctly), yet all three
        pairwise cuts fail and exact match is $0$: the cross-component logits sit just
        \emph{above} the model's own decision threshold, not far below it.
  \item \textbf{$(7,15,24)$} --- two of the three cuts succeed, cut$(2,3)$ (middle vs.\ largest)
        fails; the 24-node component also contains pairs beyond the doubled wall
        $2\cdot3^L{=}18$.
\end{itemize}

This document runs \emph{two} fine-tuning experiments off the same starting checkpoint, asking
two related but different questions.

\begin{itemize}\setlength\itemsep{0.15em}
  \item \textbf{Experiment~1 (targeted, \S\ref{sec:setup}--\S\ref{sec:verdict-a}).} Does
        fine-tuning \emph{only} the similarity read-out's two global scalars (scale, bias) ---
        everything else in the trunk frozen --- on a stream containing \emph{only these two named
        cells} fix them, without breaking the ordinary $K{=}2$ split-chains behaviour the
        checkpoint already has?
  \item \textbf{Experiment~2 (full distribution, \S\ref{sec:setup-full}--\S\ref{sec:verdict-b}).}
        Does the same two-scalar fine-tuning, given exposure to \emph{every} three-way split
        instead of just two cells, and a much longer budget, learn a genuinely \emph{general}
        three-component rule rather than overfitting the two cells above?
\end{itemize}

Both are narrower/longer than Report~IX's own fine-tuning probe
(\texttt{report/9/scratch\_prof\_experiments/scratch\_prof\_experiments.tex}, \S6), which trained
on a generic $K\in\{3,4,5,6\}$ \texttt{path\_union} stream for only $3{,}000$ steps and left
$K{\ge}3$ exact match at $0$ throughout. Fine-tuning only two scalars over a frozen trunk is
cheap, so a much longer budget costs little in either case.

\section{Experiment 1: targeted fine-tuning on the two named cells}
\label{sec:exp1}

\subsection{Setup}
\label{sec:setup}

\paragraph{Checkpoint.} The unmodified split-chains-only, $n{=}46$, seed-$1000$, similarity
read-out checkpoint from Report~IX
(\texttt{runs/report9/n46\_train/n46\_split\_chains\_roberta\_similarity\_lam0\_seed1000/last.pt}) ---
the same one behind every Report~IX \S A.3 three-way figure.

\paragraph{Fine-tuning.} Every parameter frozen except \texttt{sim\_scale}/\texttt{sim\_bias}
(2 scalars out of $\sim$6.3M). Training stream: uniformly, either the $(15,15,16)$ or the
$(7,15,24)$ three-way split-chain graph (\texttt{generate\_multi\_path\_split\_graph}), freshly
relabelled every sample. $50{,}000$ steps, batch $500$, AdamW lr $10^{-3}$
(\texttt{finetune\_readout\_threeway.py}).

\paragraph{What is measured.} Before fine-tuning (step $0$) and every $1{,}000$ steps after:
whole-graph exact match, per-component reach, the three pairwise cuts, and the overall
predicted-positive rate, on both target cells (so the effect on the \emph{held-out} cell is
visible even when only the other one dominates gradient updates); for $(7,15,24)$ specifically,
reach inside the 24-node component split into three shortest-path-distance buckets ---
within-capacity ($d\le9$), between-the-walls ($9<d\le18$), and beyond the doubled wall
($d>18$) --- the direct test of whether fine-tuning can push learning past the architectural
bound, not just recalibrate a threshold within it. Separately, \texttt{eval\_asym\_chains.py} on
the fine-tuned checkpoint reproduces the ordinary $K{=}2$ split-chains sweep (Report~IX
Table~3/\S sec:planA), compared against the identical, already-committed sweep on the
\emph{untouched} checkpoint (\texttt{runs/report9/asym\_chains\_n46/n46\_splitchains\_seed1000/}),
to check the two-scalar update did not cost anything in distribution.

\subsection{Results}
\label{sec:results}

\subsubsection{Does it learn the two target cells?}
\label{sec:res-curve}

[PENDING --- run \texttt{scripts/r10\_finetune\_threeway\_splitchains\_seed1000.sbatch}, pull
\texttt{runs/report10/}, then \texttt{python plot\_finetune\_readout\_threeway.py --tag
n46\_splitchains\_seed1000\_threeway}.]

\begin{figure}[H]
    \centering
    \figorbox{runs/report10/report10_figs/r10_finetune_threeway_curve_n46_splitchains_seed1000_threeway.png}{0.95}
    \caption{\emph{Model:} split-chains-only, $n{=}46$, seed $1000$, similarity read-out
    checkpoint of Report~IX, fine-tuned as in \S\ref{sec:setup}. \emph{Test set:} the two target
    cells, $(15,15,16)$ and $(7,15,24)$, $300$ independently relabelled graphs per evaluation.
    \emph{Metric:} whole-graph exact match, cut$(2,3)$ (the pair that fails at step $0$), and the
    overall predicted-positive rate, vs.\ fine-tuning step.}
    \label{fig:r10curve}
\end{figure}

\subsubsection{Does it learn past the doubled wall?}
\label{sec:res-beyond}

[PENDING analysis text.]

\begin{figure}[H]
    \centering
    \figorbox{runs/report10/report10_figs/r10_finetune_threeway_beyondwall_n46_splitchains_seed1000_threeway.png}{0.72}
    \caption{\emph{Model/training as in Figure~\ref{fig:r10curve}}. \emph{Test set:} cell
    $(7,15,24)$, $300$ graphs per evaluation. \emph{Metric:} pairwise reach (target connected)
    inside the 24-node component, split by shortest-path distance bucket, vs.\ fine-tuning step.}
    \label{fig:r10beyond}
\end{figure}

\subsubsection{Does it forget the $K{=}2$ split-chains behaviour?}
\label{sec:res-ownfamily}

[PENDING analysis text.]

\begin{figure}[H]
    \centering
    \figorbox{runs/report10/report10_figs/r10_finetune_threeway_ownfamily_n46_splitchains_seed1000_threeway.png}{0.72}
    \caption{\emph{Model:} the same checkpoint, before (Report~IX, untouched) vs.\ after this
    fine-tuning. \emph{Test set:} split-chain graphs at every split $a{=}1..23$, $300$ graphs per
    split. \emph{Metric:} whole-graph exact match vs.\ split $a$.}
    \label{fig:r10ownfamily}
\end{figure}

\subsubsection{How much did scale/bias actually move?}
\label{sec:res-scalebias}

[PENDING --- read off \texttt{runs/report10/finetune\_readout\_threeway/
n46\_splitchains\_seed1000\_threeway/scale\_bias\_summary.json} and comment on the magnitude of
the shift relative to the pre-fine-tuning values (comparable read-in-geometry tables in
Reports~VII--IX: scale$\approx$27--33, bias$\approx-5$ to $-8$ across every checkpoint trained on
this
distribution family so far).]

\subsubsection{Mechanistic comparison: attention and the five-stage geometry, before vs.\ after}
\label{sec:res-mechanistic}

[PENDING --- compare \texttt{runs/report9/heatmaps\_threeway/n46\_splitchains\_seed1000/} and
\texttt{runs/report9/stagewise\_threeway/n46\_splitchains\_seed1000/} (before, already committed)
against \texttt{runs/report10/heatmaps\_threeway/n46\_splitchains\_seed1000\_threeway\_post/} and
\texttt{runs/report10/stagewise\_threeway/n46\_splitchains\_seed1000\_threeway\_post/} (after,
produced by the same sbatch script): does the fine-tuned decision threshold visibly move the
final-stage ($H^{(2)}$) heatmap's cross-component blocks from ``barely positive'' to
``negative'', while leaving the within-component blocks and every earlier stage unchanged (since
the trunk itself never updates)?]

\subsection{Reading, once the data is in}
\label{sec:verdict-a}

[PENDING --- to be written after \S\ref{sec:res-curve}--\ref{sec:res-mechanistic} have real
numbers. Questions this subsection must answer plainly: (1) did it work on $(15,15,16)$; (2) did
it transfer to $(7,15,24)$ or did the two cells trade off against each other; (3) did any of the
learning happen on pairs beyond the doubled wall, or only within it; (4) was the $K{=}2$
behaviour preserved; (5) is the fix a real geometric change (visible in the mechanistic
comparison) or purely a threshold shift (scale/bias moved, nothing else did --- which, given
everything but scale/bias is frozen, is definitionally the only thing that \emph{could} move, so
this question is really about how much of the residual error the geometry itself already
carries, cf.\ Report~IX \S sec:res-readout's raw-logit reading).]

\section{Experiment 2: fine-tuning on the full $K{=}3$ split distribution}
\label{sec:exp2}

Experiment~1 asks whether the two-scalar fine-tune can fix two specific cells it is trained
directly on. It says nothing about whether that fix is a genuinely general three-component rule
or an overfit to those two cells' particular geometry. Experiment~2 asks the general-rule question
directly: fine-tune the same two scalars, from the same starting checkpoint, but on \emph{every}
three-way split instead of two, for a much longer budget.

\subsection{Setup}
\label{sec:setup-full}

\paragraph{Checkpoint.} Identical starting point to Experiment~1 (\S\ref{sec:setup}): the
unmodified split-chains-only, $n{=}46$, seed-$1000$ checkpoint.

\paragraph{Fine-tuning.} Every parameter frozen except \texttt{sim\_scale}/\texttt{sim\_bias}, as
in Experiment~1. Training stream: \texttt{generate\_path\_union\_graph(n, rng, min\_paths{=}3,
max\_paths{=}3)} --- exactly three components every sample, with both cut points (hence every
component-size combination) drawn uniformly, freshly relabelled every sample --- instead of only
the two named cells. $200{,}000$ steps (4$\times$ Experiment~1's budget), batch $500$, AdamW lr
$10^{-3}$ (\texttt{finetune\_readout\_threeway\_full.py}).

\paragraph{What is measured.} The same per-cell metrics as Experiment~1 on the \emph{same} two
named cells $(15,15,16)$/$(7,15,24)$ (for a direct, apples-to-apples comparison against
Experiment~1's curves, Figures~\ref{fig:r10curve}--\ref{fig:r10beyond}), \emph{plus} an aggregate
exact-match/pairwise-accuracy number over freshly sampled, generic three-component graphs at every
evaluation (the same style of number Report~IX's own generic-$K{\ge}3$ fine-tuning probe
reported). The own-family $K{=}2$ sweep and the mechanistic battery (attention, five-stage
geometry) are reproduced exactly as in Experiment~1, into the \texttt{\_full3way}-tagged output
directories.

\subsection{Results}
\label{sec:results-full}

\subsubsection{Does exposure to the full distribution generalise, where the targeted fine-tune
could only overfit two cells?}
\label{sec:res-full-aggregate}

[PENDING --- run \texttt{scripts/r10\_finetune\_full3way\_splitchains\_seed1000.sbatch}, pull
\texttt{runs/report10/}, then \texttt{python plot\_finetune\_readout\_threeway.py --tag
n46\_splitchains\_seed1000\_full3way --run\_subdir finetune\_readout\_full3way}.]

\begin{figure}[H]
    \centering
    \figorbox{runs/report10/report10_figs/r10_finetune_full3way_aggregate_n46_splitchains_seed1000_full3way.png}{0.72}
    \caption{\emph{Model:} split-chains-only, $n{=}46$, seed $1000$, similarity read-out
    checkpoint of Report~IX, fine-tuned as in \S\ref{sec:setup-full}. \emph{Test set:} freshly
    sampled, generic three-component \texttt{path\_union} graphs, $300$ per evaluation.
    \emph{Metric:} whole-graph exact match and pairwise accuracy, vs.\ fine-tuning step.}
    \label{fig:r10fullaggregate}
\end{figure}

\subsubsection{Same two named cells as Experiment 1 --- does the broader stream do better,
worse, or the same?}
\label{sec:res-full-named}

[PENDING --- regenerate Figures~\ref{fig:r10curve}--\ref{fig:r10beyond}-equivalent plots with
\texttt{--tag n46\_splitchains\_seed1000\_full3way --run\_subdir finetune\_readout\_full3way}
and compare the resulting curves directly against Experiment~1's, side by side.]

\subsubsection{Own-family retention and scale/bias movement}
\label{sec:res-full-ownfamily}

[PENDING --- same construction as \S\ref{sec:res-ownfamily}/\S\ref{sec:res-scalebias}, on the
\texttt{\_full3way}-tagged outputs.]

\subsection{Reading, once the data is in}
\label{sec:verdict-b}

[PENDING. Central question: does the aggregate K=3 exact match (Figure~\ref{fig:r10fullaggregate})
move meaningfully above $0$, i.e.\ does read-out-only fine-tuning ever produce a genuinely general
three-component rule given enough coverage and steps --- or does it plateau near $0$ the way
Report~IX's shorter, mixed-$K$ probe did, just confirming at higher confidence (longer run, purer
$K{=}3$ stream) that the representational gap is in the frozen trunk, not the read-out threshold?
Compare directly against Experiment~1: if the full-distribution run does \emph{worse} on the two
named cells than the targeted run, that is itself informative (spreading two scalars' worth of
capacity over the whole $K{=}3$ space costs the two hardest cells something the narrow fine-tune
did not have to pay).]

\section{Combined verdict}
\label{sec:verdict}

[PENDING --- write only after both \S\ref{sec:verdict-a} and \S\ref{sec:verdict-b} are filled in.
This is the section whose conclusion actually matters for the paper's \S5.3: state plainly whether
\emph{any} amount of read-out-only fine-tuning (targeted or general, $50$k or $200$k steps) gets a
three-component graph right, and if so under which condition specifically, in one or two
sentences Edoardo can lift almost verbatim into a conversation with the professor about what
\S5.3 should claim.]

\end{document}
